\documentclass[11pt]{article}
\usepackage[margin=1.1in]{geometry}
\usepackage{amsmath,amssymb,amsthm}
\usepackage{hyperref}

\newtheorem{theorem}{Theorem}[section]
\newtheorem{lemma}[theorem]{Lemma}
\newtheorem{corollary}[theorem]{Corollary}
\theoremstyle{definition}
\newtheorem{definition}[theorem]{Definition}
\newtheorem{conjecture}[theorem]{Conjecture}
\theoremstyle{remark}
\newtheorem*{remark}{Remark}

\title{A Complete Proof of Conjecture \textup{TLMC-00000001206}:\\ the Nim-Sum Period Law}
\author{Solution submission}
\date{September 12, 2026}

\begin{document}
\maketitle

\begin{abstract}
We prove Conjecture \textup{TLMC-00000001206} of the Last Math Competition, which asserts that the period of the pointwise nim-sum of two periodic Grundy sequences divides the least common multiple of the two periods. We give an elementary, self-contained, unconditional proof, together with a machine-verified Lean~4 formalization (file \texttt{NimSumPeriod.lean}, no \texttt{sorry}). In fact we prove a stronger statement: the conclusion holds for \emph{arbitrary} $\mathbb{N}$-valued periodic sequences under \emph{any} pointwise binary operation; the conjecture for Grundy sequences and the nim-sum follows a fortiori.
\end{abstract}

\section{The conjecture}

The conjecture as recorded in \texttt{conjectures/00000001206.md} reads:

\begin{conjecture}[TLMC-00000001206; nim-sum period]\label{conj:main}
\emph{(English)} Nim-sum closure of Sprague--Grundy functions: the period of the pointwise nim-sum of two periodic Grundy sequences divides the lcm of the two periods.
\par\smallskip
\emph{(Chinese)} Sprague--Grundy 函数的 nim-和封闭：两周期 Grundy 序列的逐点 nim-和的周期为两周期之 lcm 的因子。
\end{conjecture}

\section{Definitions and notation}

Throughout, all sequences are functions $g \colon \mathbb{N} \to \mathbb{N}$, where $\mathbb{N} = \{0, 1, 2, \dots\}$.

\begin{definition}[Period]\label{def:period}
A positive integer $p \ge 1$ is a \emph{period} of a sequence $g \colon \mathbb{N} \to \mathbb{N}$ if
\[
g(n + p) = g(n) \qquad \text{for all } n \in \mathbb{N}.
\]
A sequence is \emph{periodic} if it has at least one period.
\end{definition}

\begin{definition}[Nim-sum]\label{def:nimsum}
The \emph{nim-sum} (bitwise exclusive or) of two nonnegative integers is written $x \oplus y$; explicitly, if $x = \sum_i x_i 2^i$ and $y = \sum_i y_i 2^i$ with digits $x_i, y_i \in \{0,1\}$, then $x \oplus y = \sum_i (x_i + y_i \bmod 2)\, 2^i$. The \emph{pointwise nim-sum} of two sequences $g_1, g_2$ is the sequence $n \mapsto g_1(n) \oplus g_2(n)$.
\end{definition}

\begin{definition}[Minimal period]\label{def:minperiod}
The \emph{minimal period} (or \emph{the} period) $\pi(g)$ of a periodic sequence $g$ is the least positive period:
\[
\pi(g) \;=\; \min\,\{\, p \ge 1 :\ p \text{ is a period of } g \,\}.
\]
The minimum exists because the set is a nonempty set of positive integers.
\end{definition}

\begin{definition}[Grundy sequence]\label{def:grundy}
A \emph{Grundy sequence} is the Sprague--Grundy function $g$ of an impartial combinatorial game (with positions indexed by $\mathbb{N}$): $g(n) = \operatorname{mex}\{\, g(m) : m \in \mathrm{opt}(n) \,\}$, where $\operatorname{mex}$ is the least excluded nonnegative integer. In particular, every Grundy sequence is an $\mathbb{N}$-valued sequence.
\end{definition}

\section{The theorem and its complete proof}

\begin{theorem}[Nim-sum period law]\label{thm:main}
Let $g_1, g_2 \colon \mathbb{N} \to \mathbb{N}$ be periodic sequences with periods $p_1, p_2 \ge 1$, respectively. Then the pointwise nim-sum $g = (n \mapsto g_1(n) \oplus g_2(n))$ is periodic, and its minimal period divides $\operatorname{lcm}(p_1, p_2)$:
\[
\pi(g) \;\mid\; \operatorname{lcm}(p_1, p_2).
\]
\end{theorem}

\begin{proof}
The proof uses only the fact that the nim-sum is a \emph{function} of the two values $g_1(n)$ and $g_2(n)$; we isolate three elementary steps.

\emph{Step 1 (multiples of a period are periods).} Let $f$ be a sequence with a period $p \ge 1$. We show by induction on $k$ that $k \cdot p$ is a period of $f$ for every $k \in \mathbb{N}$. For $k = 0$ this says $f(n + 0) = f(n)$, which is trivial. Assume $k \cdot p$ is a period. Then for every $n$,
\[
f\bigl(n + (k+1)p\bigr) \;=\; f\bigl((n + k p) + p\bigr) \;\overset{(\ast)}{=}\; f(n + k p) \;=\; f(n),
\]
where $(\ast)$ holds because $p$ is a period (applied at the index $n + kp$), and the last step is the induction hypothesis. Hence $(k+1) \cdot p$ is a period, completing the induction.

\emph{Step 2 (the lcm is a period of the nim-sum).} Write $L = \operatorname{lcm}(p_1, p_2)$. Since $p_1 \mid L$, there is $a \in \mathbb{N}$ with $L = a\,p_1$; by Step 1 (with $k = a$), $g_1(n + L) = g_1(n)$ for all $n$. Similarly $p_2 \mid L$ gives $b \in \mathbb{N}$ with $L = b\,p_2$ and $g_2(n + L) = g_2(n)$ for all $n$. Consequently, for every $n$,
\[
g(n + L) \;=\; g_1(n + L) \oplus g_2(n + L) \;=\; g_1(n) \oplus g_2(n) \;=\; g(n),
\]
so $L$ is a period of $g$. In particular $g$ is periodic, and by Definition~\ref{def:minperiod} its minimal period $\pi = \pi(g)$ exists and satisfies $1 \le \pi \le L$.

\emph{Step 3 (the minimal period divides every period).} Let $q \ge 1$ be any period of $f$, and let $\pi$ be the minimal period of $f$; we claim $\pi \mid q$. Perform Euclidean division: $q = c\,\pi + r$ with $c \in \mathbb{N}$ and $0 \le r < \pi$. By Step 1, $c\,\pi$ is a period of $f$, so for every $n$,
\[
f(n + r + c\,\pi) = f(n + r).
\]
On the other hand $n + r + c\,\pi = n + q$, and since $q$ is a period, $f(n + q) = f(n)$. Combining,
\[
f(n + r) \;=\; f(n + q) \;=\; f(n) \qquad \text{for all } n.
\]
Thus $r$ is a period of $f$. If $r \ge 1$, then $r$ is a positive period with $r < \pi$, contradicting the minimality of $\pi$. Hence $r = 0$, i.e.\ $\pi \mid q$, proving the claim.

\emph{Conclusion.} Apply Step 3 to $f = g$, $\pi = \pi(g)$, and the period $q = L$ supplied by Step 2: $\pi(g) \mid L = \operatorname{lcm}(p_1, p_2)$, which is the assertion of the theorem.
\end{proof}

\section{Derivation of Conjecture~\ref{conj:main}}

\begin{corollary}\label{cor:conj}
Conjecture~\ref{conj:main} holds.
\end{corollary}

\begin{proof}
Let $g_1, g_2$ be periodic Grundy sequences with (minimal) periods $p_1, p_2$. Grundy sequences are $\mathbb{N}$-valued (Definition~\ref{def:grundy}), and the pointwise nim-sum of two $\mathbb{N}$-valued sequences is $\mathbb{N}$-valued. Theorem~\ref{thm:main} applies to arbitrary $\mathbb{N}$-valued periodic sequences; hence the minimal period of $n \mapsto g_1(n) \oplus g_2(n)$ divides $\operatorname{lcm}(p_1, p_2)$, which is precisely the assertion of Conjecture~\ref{conj:main}.
\end{proof}

\begin{remark}[Strength of the result]
The proof of Theorem~\ref{thm:main} uses nothing about the nim-sum beyond its being a well-defined binary operation on the value set: Steps 1--3 are statements about periods alone. The identical argument proves: \emph{for any binary operation $\star \colon \mathbb{N} \times \mathbb{N} \to \mathbb{N}$ (pointwise addition, multiplication, maximum, bitwise and/or/xor, \dots), the minimal period of $n \mapsto g_1(n) \star g_2(n)$ divides $\operatorname{lcm}(p_1, p_2)$}. The nim-sum case of the conjecture is one instance. The Lean formalization is carried out for a completely arbitrary binary operation, making this generality machine-checked.
\end{remark}

\begin{remark}[Sharpness]
Both divisibility conclusions are optimal in general: taking $g_2$ constant, the nim-sum $g_1 \oplus g_2$ need not have any period smaller than the lcm of the individual minimal periods, so no strengthening of the divisibility (e.g.\ to a fixed proper factor of $\operatorname{lcm}(p_1,p_2)$) is possible without further hypotheses.
\end{remark}

\section{Lean 4 formalization}

The file \texttt{NimSumPeriod.lean} (toolchain \texttt{leanprover/lean4:v4.31.0}, Mathlib \texttt{v4.31.0}) contains a complete machine-checked proof:
\begin{itemize}
\item \texttt{IsPeriod f p} --- Definition~\ref{def:period};
\item \texttt{minPeriod f} --- Definition~\ref{def:minperiod}, via least positive period;
\item \texttt{IsPeriod.mul}, \texttt{minPeriod\_dvd} --- Steps 1 and 3;
\item \texttt{pointwise\_period} --- Step 2 and Theorem~\ref{thm:main}, proved for an \emph{arbitrary} binary operation \texttt{op} in place of \texttt{Nat.xor};
\item \texttt{nimsum\_period} --- Theorem~\ref{thm:main} specialized to the nim-sum \texttt{(\^\^\^)} = \texttt{Nat.xor};
\item \texttt{conjecture1206\_true} --- the faithful statement of Conjecture~\ref{conj:main} for $\mathbb{N}$-valued sequences, proved.
\end{itemize}
The file compiles with no errors and contains no \texttt{sorry}.

\section{Conclusion}

Conjecture \textup{TLMC-00000001206} is \textbf{proved}: for any two $\mathbb{N}$-valued periodic sequences with periods $p_1, p_2 \ge 1$, the minimal period of their pointwise nim-sum divides $\operatorname{lcm}(p_1, p_2)$ (Theorem~\ref{thm:main}, Corollary~\ref{cor:conj}). The proof is elementary, unconditional, and fully formalized in Lean~4.

\end{document}
